\subsubsection{Tendência de Engajamento}
A Tendência de Engajamento \textit{Engagement Trend} analisa a evolução temporal do foco dos alunos ao longo da duração total da aula, dividindo a sessão em quatro segmentos temporais iguais (quartis). Esta métrica permite identificar padrões de variação no engajamento, como desgaste da atenção, picos de interesse ou momentos críticos de distração.

O cálculo é realizado pela função \textbf{calculate\_engagement\_trend(traces, lecture\_id)} que implementa a seguinte abordagem:

\begin{equation}
\text{Engajamento}(Q_i) = \frac{\sum \text{tempo focado em } Q_i}{\sum \text{tempo total em } Q_i} \quad \text{para } i \in \{1,2,3,4\}
\end{equation}

% \lstinputlisting[
%     language=python,
%     caption={Cálculo da Tendência de Engajamento},
%     firstline=11,
%     lastline=169,
%     label={lst:calculate_engagement_trend}
% ]{scripts/metricsCalc/time.py}

\textbf{Etapas do cálculo}:
\begin{enumerate}
    \item \textbf{Recuperação de metadados}: Obtém data, horário de início e fim reais da aula
    \item \textbf{Construção temporal}: Cria objetos \textit{datetime} com data e hora combinadas
    \item \textbf{Definição de quartis}: Divide a duração total em quatro intervalos iguais
    \item \textbf{Mapeamento temporal}: Para cada intervalo entre \textit{traces} consecutivos:
    \begin{itemize}
        \item Calcula tempo decorrido desde o início da aula
        \item Atribui ao quartil correspondente baseado no timestamp
        \item Soma ao tempo total do quartil
        \item Se estava focado na aula, soma ao tempo focado do quartil
    \end{itemize}
    \item \textbf{Cálculo de proporções}: Para cada quartil, calcula razão entre tempo focado e tempo total
\end{enumerate}

\textbf{Definição dos quartis}:
\begin{itemize}
    \item \textbf{Q1 (Quartil 1)}: Primeiro 25\% da duração da aula (0--25\%)
    \item \textbf{Q2 (Quartil 2)}: Segundo 25\% da duração (25--50\%)
    \item \textbf{Q3 (Quartil 3)}: Terceiro 25\% da duração (50--75\%)
    \item \textbf{Q4 (Quartil 4)}: Último 25\% da duração (75--100\%)
\end{itemize}


\subsubsection{Pico de Engajamento}
O Pico de Engajamento Temporal (\textit{Peak Engagement Time}) identifica o período de 5 minutos durante a aula que concentrou a maior quantidade de tempo focado pelos alunos. Esta métrica permite localizar os momentos de máxima atenção coletiva, fornecendo insights sobre quais segmentos da aula foram mais efetivos em capturar e manter o engajamento dos estudantes.

O cálculo é realizado pela função \textbf{calculate\_peak\_engagement\_time(traces, lecture\_id)} que implementa a seguinte abordagem:

\begin{equation}
\text{Pico de Engajamento} = \arg\max_{t \in T} \sum_{a \in A} \text{tempo\_focado}_a(t)
\end{equation}

\textbf{onde:}
\begin{itemize}
    \item $T$ é o conjunto de intervalos de 5 minutos durante a aula
    \item $A$ é o conjunto de todos os alunos
    \item $\text{tempo\_focado}_a(t)$ é o tempo que o aluno $a$ esteve focado durante o intervalo $t$
\end{itemize}

% \lstinputlisting[
%     language=python,
%     caption={Cálculo do Pico de Engajamento},
%     firstline=201,
%     lastline=255,
%     label={lst:calculate_peak_engagement_time}
% ]{scripts/metricsCalc/time.py}

\textbf{Etapas do cálculo}:
\begin{enumerate}
    \item \textbf{Recuperação temporal}: Obtém horários reais de início e fim da aula
    \item \textbf{Agrupamento em intervalos}: Divide a aula em segmentos de 5 minutos
    \item \textbf{Processamento incremental}: Para cada intervalo entre \textit{traces}:
    \begin{itemize}
        \item Determina o intervalo de 5 minutos correspondente ao timestamp
        \item Soma o tempo de foco (quando na aba da aula) ao intervalo
        \item Acumula tempo focado de todos os alunos
    \end{itemize}
    \item \textbf{Identificação do pico}: Localiza o intervalo com maior soma acumulada de tempo focado
\end{enumerate}

\subsubsection{Ponto de Queda de Engajamento}
O Ponto de Queda de Engajamento (\textit{Dropoff Point}) identifica o momento temporal em que o nível de foco dos alunos cai para 50\% do pico máximo de engajamento observado durante a aula e permanece abaixo deste limiar de forma sustentada. Esta métrica é crucial para detectar quando a atenção coletiva diminui significativamente, indicando possíveis pontos de fadiga cognitiva ou perda de interesse no conteúdo.

O cálculo é realizado pela função \textbf{calculate\_dropoff\_point(traces, lecture\_id)} que implementa a seguinte lógica:

\begin{equation}
\text{Dropoff Point} = \min\{t \mid \text{Engajamento}(t) < 0.5 \times \text{Pico} \land \forall s \in [t, t+15]: \text{Engajamento}(s) < 0.5 \times \text{Pico}\}
\end{equation}

\textbf{onde:}
\begin{itemize}
    \item $\text{Pico} = \max_{t \in T} \text{Engajamento}(t)$ é o nível máximo de engajamento
    \item $\text{Engajamento}(t)$ é a proporção de foco no intervalo temporal $t$
    \item O limiar é definido como 50\% do pico máximo
\end{itemize}

% \lstinputlisting[
%     language=python,
%     caption={Cálculo do Declive no Engajamento},
%     firstline=257,
%     lastline=417,
%     label={lst:calculate_dropoff_point}
% ]{scripts/metricsCalc/time.py}

\textbf{Etapas do cálculo}:
\begin{enumerate}
    \item \textbf{Processamento temporal}: Converte timestamps para tempo real desde início da aula
    \item \textbf{Agrupamento em intervalos}: Divide em segmentos de 5 minutos
    \item \textbf{Cálculo de engajamento}: Para cada intervalo, calcula proporção de tempo focado
    \item \textbf{Identificação do pico}: Determina nível máximo de engajamento
    \item \textbf{Definição do limiar}: Estabelece limite como 50\% do pico máximo
    \item \textbf{Detecção de queda}: Encontra primeiro intervalo abaixo do limiar com sustentação temporal
    \item \textbf{Verificação de sustentação}: Confirma que queda persiste por 10-15 minutos
\end{enumerate}

\textbf{Critérios de detecção}:
\begin{itemize}
    \item \textbf{Limiar adaptativo}: 50\% do pico máximo (não valor absoluto fixo)
    \item \textbf{Sustentação temporal}: Queda deve persistir por 10-15 minutos
    \item \textbf{Intervalos consecutivos}: Próximos 2-3 intervalos (10-15min) também abaixo do limiar
    \item \textbf{Foco no primeiro ocorrência}: Identifica ponto inicial da queda significativa
\end{itemize}